\documentclass{vkr} % твой класс документа


\usepackage{array}
\usepackage{tabularx}
\usepackage{longtable}
\usepackage{multirow}
\usepackage{booktabs}
\usepackage{caption}
\usepackage{float}


% ------------------- Языки -------------------
\usepackage[english, russian]{babel} % русский + английский
\usepackage[utf8]{inputenc} % кодировка
\usepackage[T2A]{fontenc} % шрифты для русского

% ------------------- Математика -------------------
\usepackage{amsmath}
\usepackage{amssymb}

% ------------------- Рисунки -------------------
\usepackage{graphicx}
\graphicspath{{images/}} % папка с изображениями
\usepackage{float} % позиционирование рисунков [H]
\usepackage{pdflscape} % альбомные страницы

% ------------------- Ссылки -------------------
\usepackage[hidelinks]{hyperref}

% ------------------- Таблицы -------------------
\usepackage{xltabular}
\usepackage{array}
\usepackage{makecell}
\renewcommand\theadfont{} % шрифт заголовков таблиц
\newcolumntype{T}{>{\centering\arraybackslash}X} % центрированная колонка
\newcolumntype{R}{>{\raggedleft\arraybackslash}X} % выравнивание по правому краю
\newcolumntype{C}[1]{>{\centering\arraybackslash}m{#1}} % фиксированная ширина, центр
\newcolumntype{r}[1]{>{\raggedleft\arraybackslash}p{#1}} % фиксированная ширина, вправо
\newcommand{\centrow}{\centering\arraybackslash} % центрирование ячейки
\newcommand{\finishhead}{\endhead\hline\endlastfoot} % для xltabular
\usepackage[tableposition=top]{caption} % подписи таблиц
\captionsetup[table]{font={stretch=1.2}, skip=0pt, belowskip=0pt, aboveskip=8pt, singlelinecheck=off}

% ------------------- Листинги -------------------
\usepackage{listings}
\usepackage{xcolor}

% Цвета для Python
\definecolor{codegreen}{rgb}{0,0.6,0}
\definecolor{codegray}{rgb}{0.5,0.5,0.5}
\definecolor{codepurple}{rgb}{0.58,0,0.82}

\lstset{
	language=Python,
	numbers=left,
	numberstyle=\tiny\color{codegray},
	stepnumber=1,
	numbersep=5pt,
	commentstyle=\color{codegreen},
	keywordstyle=\color{magenta},
	stringstyle=\color{codepurple},
	basicstyle=\ttfamily\small,
	backgroundcolor=\color{white},
	showspaces=false,
	showstringspaces=false,
	showtabs=false,
	frame=single,
	tabsize=2,
	captionpos=t,
	breaklines=true,
	breakatwhitespace=false,
	escapeinside={\%*}{*)}
}

% ------------------- Отступы списков -------------------
\usepackage{enumitem}
\setlist{nolistsep, wide=\parindent, itemindent=*}

% ------------------- Стиль списка источников -------------------
\addto{\captionsrussian}{\renewcommand{\refname}{СПИСОК ИСПОЛЬЗОВАННЫХ ИСТОЧНИКОВ}}

% ------------------- Полезные настройки -------------------
\usepackage{setspace} % межстрочный интервал
\setstretch{1.2}
